\chapter{引 言}

\section{课题背景}

操作系统课程需要通过可运行的内核帮助学生理解抽象机制。仅依赖概念讲授时，进程调度、地址空间、文件系统和设备驱动之间的关系不易呈现\cite{you2000unix}。Unix V6++（Unix V6++ Tongji's Edition for Education）是同济大学陈闳中教授及多届研究生团队开发的教学操作系统。该系统以 Unix Version 6 为基础，在 IA32 架构上使用 C++ 重新实现，代码规模较小，覆盖进程管理、内存管理、文件系统和设备管理等基本模块，因而适合作为课程实验和源码阅读对象\cite{chen2013unixv6pp}。

原系统的教学价值并不意味着其实现可以长期不变。IA32 架构已不再是主流教学与实验平台，原系统中与 x86 分页、中断控制器和 ATA 设备绑定的代码增加了迁移成本。C++ 实现虽然通过类封装组织了部分内核对象，但底层仍大量依赖裸指针、全局状态和手工引用计数；在文件对象、inode、缓冲区和页对象等共享资源上，生命周期关系主要由程序员约定维护。该设计便于展示传统 Unix 实现方式，但在后续扩展和调试中容易形成隐式依赖。

RISC-V 和 Rust 为教学操作系统的重构提供了两个现实选择。RISC-V 指令集结构相对规整，特权级、异常和分页机制文档公开，适合用于解释操作系统与硬件之间的接口\cite{waterman2017riscv}。Rust 则通过所有权、借用检查和 RAII 风格析构机制，将一部分资源管理约束前移到类型系统中\cite{klabnik2022rust}。MIT xv6 已完成 RISC-V 版本迁移\cite{kaashoek2021xv6}，rCore 和 Writing an OS in Rust 展示了 Rust 在教学内核中的应用方式\cite{chen2023rcore,oppermannWritingOS}。工业界也开始在操作系统内核或内核组件中引入 Rust 技术栈\cite{antgroup2023asterinas}。在这一背景下，对 Unix V6++ 的内存管理和文件系统进行 Rust 重构，并探索 RISC-V 迁移，能够在保留原系统课程价值的同时降低后续维护和扩展成本。

\section{相关研究现状}

教学操作系统通常在完整性和可读性之间取舍。Unix V6 以较小规模呈现了早期 Unix 内核的核心设计，Lions 对 Unix 第六版的注释长期用于操作系统教学\cite{lions1996unix}。xv6 延续了这一思路，通过更现代的代码组织重写 Unix V6，并在 2019 年后将主要教学平台转向 RISC-V\cite{kaashoek2021xv6}。这类系统的共同特点是功能有限，但模块边界相对清晰，便于围绕进程、内存、文件和设备展开实验。

Rust 教学内核则更关注类型系统对系统软件的约束能力。rCore 使用 Rust 从零实现操作系统，并将所有权、生命周期和错误处理纳入教学内容\cite{chen2023rcore}。Writing an OS in Rust 通过逐步构建内核的方式说明了 Rust 在裸机环境中的基本用法\cite{oppermannWritingOS}。这些工作表明，Rust 可以用于组织页表、堆分配器和驱动等底层模块，但 unsafe 边界仍然无法完全消除。对教学系统而言，关键不在于宣称完全安全，而在于将裸指针和硬件相关操作限制在较小范围，并在上层提供语义明确的接口。

Unix V6++ 与上述系统相比具有明确的本地课程背景。它既保留 Unix V6 的教学结构，又承载同济大学操作系统课程的实验积累。本课题的研究重点并非重新设计一个全新的教学内核，而是在原系统结构上比较 C++、Rust IA32 和 Rust RISC-V 三个版本的差异，尤其关注内存管理和文件系统在资源生命周期、错误处理和跨架构复用方面的变化。

\section{研究目标}

本课题围绕 Unix V6++ 的内存管理和文件系统展开。内存管理部分使用 Rust 重新组织物理页、页分配器、内核堆和交换区管理，保留原有页式管理和教学用地址空间布局，同时引入 Buddy 与 Slab 分层分配、类型化地址表示和 RAII 风格资源回收。文件系统部分保留超级块、inode、目录项、打开文件表和缓冲区缓存等 Unix V6 风格结构，重点重构对象引用关系、错误传播路径和文件描述符管理逻辑。

在体系结构迁移方面，本课题将 Rust 版本适配到 RISC-V 64 位平台，完成 Sv39 分页、内核与用户空间映射、UART 串口和 VirtIO 块设备等基础路径。论文通过对 main 分支原始 C++ 版本、refactor/x86 分支和 refactor/riscv 分支的比较，分析 Rust 重构和 RISC-V 迁移对模块边界、可维护性和运行表现的影响。最终目标是在不削弱 Unix V6++ 教学连续性的前提下，为后续课程实验和系统扩展提供一套结构更清晰的实现基础。
